\documentclass[11pt]{article}
\usepackage{amsmath, amssymb, amsthm}

\title{Triangle Mutiny: Activation Algebra v2.7}
\author{Borealis S. Hedling}
\date{17 October 2026}

\begin{document}
\maketitle

\section*{Overview}

Activation Algebra v2.7 provides algebraic structure for activation modeling
without performing activation. This layer consumes activation semantics (v2.6),
activation modeling (v2.5), activation representation (v2.4), activation
simulation (v2.3), activation transitions (v2.2), activation surfaces (v2.1),
and activation logic (v2.0).

The goal of this layer is to define algebraic operators, relations, bundles,
and invariants that describe activation behavior in a deterministic and
non-activating manner.

\section*{Activation Models and States}

Let $M$ denote the set of activation models defined in v2.5:

\[
M = \{ (s_{\text{stable}}, s_{\text{spectral}}, s_{\text{collapse}},
        s_{\text{drift}}, s_{\text{modal}},
        v_{\text{sim}}, v_{\text{repr}} ) \}
\]

Let $S$ denote the set of symbolic activation states defined in v2.4:

\[
S = \{ (s_{\text{stable}}, s_{\text{spectral}}, s_{\text{collapse}},
        s_{\text{drift}}, s_{\text{modal}} ) \}
\]

\section*{Activation Operator}

We define an algebraic activation operator:

\[
\mathcal{A} : M \times S \to M
\]

such that:

\[
\mathcal{A}(M, S) =
  (s_{\text{stable}} \land S_{\text{stable}},
   s_{\text{spectral}} \land S_{\text{spectral}},
   s_{\text{collapse}} \land S_{\text{collapse}},
   s_{\text{drift}} \land S_{\text{drift}},
   s_{\text{modal}} \land S_{\text{modal}},
   v_{\text{sim}}, v_{\text{repr}})
\]

This operator is deterministic and non-activating.

\section*{Activation Relation}

We define an activation relation:

\[
R \subseteq M \times M
\]

such that:

\[
(M_1, M_2) \in R
\quad \text{iff} \quad
M_1.\text{modelValid} \land M_2.\text{modelValid}
\]

This relation captures semantic and geometric compatibility.

\section*{Activation Bundle}

We define an activation bundle:

\[
\mathcal{B} = (\mathcal{A}, R)
\]

representing the algebraic structure of activation.

\section*{Algebraic Invariants}

We define algebraic invariants:

\[
I_{\text{geom}}(M) =
  M_{\text{stable}} \land M_{\text{spectral}} \land M_{\text{collapse}}
\]

\[
I_{\text{prov}}(M) =
  M_{\text{simulationValid}} \land M_{\text{representationValid}}
\]

\[
I_{\text{sem}}(M) =
  M_{\text{modelValid}}
\]

\section*{Composition}

We define composition of activation operators:

\[
\mathcal{A}(M_1, S_1) \circ \mathcal{A}(M_2, S_2)
  = \mathcal{A}(M_1, S_1) \text{ followed by } \mathcal{A}(M_2, S_2)
\]

Composition is associative:

\[
(\mathcal{A}(M_1, S_1) \circ \mathcal{A}(M_2, S_2)) \circ
  \mathcal{A}(M_3, S_3)
=
\mathcal{A}(M_1, S_1) \circ
  (\mathcal{A}(M_2, S_2) \circ \mathcal{A}(M_3, S_3))
\]

\section*{Identity Element}

We define an identity activation state:

\[
e = ( \text{true}, \text{true}, \text{true},
      \text{true}, \text{true} )
\]

such that:

\[
\mathcal{A}(M, e) = M
\]

\section*{Ordering}

We define a partial order on activation models:

\[
M_1 \preceq M_2
\quad \text{iff} \quad
M_1.\text{stable} \Rightarrow M_2.\text{stable},
\quad
M_1.\text{spectral} \Rightarrow M_2.\text{spectral},
\quad
M_1.\text{collapse} \Rightarrow M_2.\text{collapse}
\]

This ordering reflects geometric clearance.

\section*{Algebraic Kernel}

The algebraic kernel is defined as:

\[
K_{\text{alg}}(M) =
  I_{\text{geom}}(M) \land I_{\text{prov}}(M) \land I_{\text{sem}}(M)
\]

\section*{Summary}

Activation Algebra v2.7 provides:

\begin{itemize}
\item algebraic activation operator,
\item activation relation,
\item activation bundle,
\item algebraic invariants,
\item operator composition,
\item identity element,
\item partial ordering,
\item algebraic kernel.
\end{itemize}

This layer prepares the system for Activation Formalism v2.8, where activation
will be axiomatized.

\end{document}
